% ============================================================================
%  Shared preamble for the combined CS221 mini-project report.
%  Union of the macros used by the seven official CS221 templates (see
%  report/templates/), adapted so the whole report compiles with plain
%  pdflatex/latexmk: no minted, so no -shell-escape is needed.
% ============================================================================

% ---- page & fonts -----------------------------------------------------------
\usepackage[margin=1in]{geometry}
\usepackage[T1]{fontenc}
\usepackage[utf8]{inputenc}
\usepackage{lmodern}
\usepackage{microtype}
\usepackage{fancyhdr}
\setlength{\parindent}{0pt}
\setlength{\parskip}{0.4em}
\setlength{\headheight}{14pt}

% ---- maths ------------------------------------------------------------------
\usepackage{amsmath,amssymb,amsthm,mathtools,mathrsfs,bm}
\DeclareMathOperator*{\argmax}{arg\,max}
\DeclareMathOperator*{\argmin}{arg\,min}
\DeclareMathOperator{\E}{\mathbb{E}}
\DeclareMathOperator{\Var}{Var}
\DeclareMathOperator{\Cov}{Cov}
\DeclareMathOperator{\sign}{sign}
\newcommand{\R}{\mathbb{R}}
\renewcommand{\P}{\mathbb{P}}
\newcommand{\bvec}[1]{\mathbf{#1}}
\newcommand{\norm}[2][2]{\| #2\|_{#1}}
\newcommand{\indic}[1]{\mathbf{1}[#1]}
% sentiment (Assignment 2, Problem 4)
\newcommand{\TrainLossNew}{\text{TrainLoss}_{\text{new}}}
\newcommand{\TrainLossAvg}{\text{TrainLoss}_{\text{avg}}}
\newcommand{\TrainLossMax}{\text{TrainLoss}_{\max}}
\newcommand{\TrainLossG}{\text{TrainLoss}_g}

\theoremstyle{definition}
\newtheorem*{lemma}{Lemma}
\newtheorem*{claim}{Claim}
\newtheoremstyle{case}{}{}{}{}{}{:}{ }{}
\theoremstyle{case}
\newtheorem{case}{Case}

% ---- tables, figures, diagrams ---------------------------------------------
\usepackage{graphicx}
\graphicspath{{figures/}}
\usepackage{booktabs,tabularx,multirow,array}
\usepackage{float}
\usepackage{caption,subcaption}
\usepackage{tikz}
\usetikzlibrary{arrows.meta,positioning,shapes,fit,calc}
% Bayesian-network node styles (Assignment 6): \node[latent] / \node[obs]
\tikzset{
  latent/.style={circle,draw,minimum size=11mm,inner sep=1pt},
  obs/.style={latent,fill=gray!30},
  edge/.style={-{Latex[length=2mm]}}
}

% ---- lists, code, colour, links --------------------------------------------
\usepackage{enumitem}
\usepackage[table]{xcolor}
\usepackage{listings}
\lstset{basicstyle=\ttfamily\small, columns=fullflexible, breaklines=true,
        frame=single, keywordstyle=\color{blue!60!black}, commentstyle=\color{green!40!black},
        language=Python, showstringspaces=false}
\usepackage{framed}
\definecolor{shadecolor}{gray}{0.93}
\usepackage{csquotes}

% ---- section styling: "Assignment N: ..." / "Problem k: ..." ---------------
\usepackage{titlesec}
\titleformat{\section}{\newpage\LARGE\bfseries}{Assignment \thesection:}{0.5em}{}
\titleformat{\subsection}{\Large\bfseries}{Problem \thesubsection:}{0.5em}{}
\renewcommand{\thesubsection}{\arabic{subsection}}
\usepackage{titletoc}
\titlecontents{section}[0pt]{\addvspace{0.6em}\bfseries}
  {Assignment \thecontentslabel:\quad}{}{\hfill\contentspage}
\titlecontents{subsection}[1.5em]{}
  {Problem \thecontentslabel:\quad}{}{\titlerule*[0.6pc]{.}\contentspage}
\setcounter{tocdepth}{2}

% ---- question / answer macros ----------------------------------------------
% \problempart{1a}{Short title}   heading for one lettered part of a problem
% \question{...}                  the prompt, restated so the PDF is self-contained
% \expect{...}                    what the course staff expect in the answer
% \begin{solution} ... \end{solution}   OUR answer (shaded, breaks across pages)
% \codepart{N.x}                  pointer to the notebook for programming parts
% \todo{...}                      red placeholder; must be gone before submission
\newcommand{\problempart}[2]{\subsubsection*{(#1)\ #2}}
\newcommand{\question}[1]{#1\par}
\newcommand{\note}[1]{\medskip\noindent\textbf{Note:} #1\par}
\newcommand{\hint}[1]{\medskip\noindent\textit{Hint:} #1\par}
\newcommand{\expect}[1]{\medskip\noindent\fbox{\parbox{0.97\linewidth}{\textbf{What we expect:} #1}}\par}
\newenvironment{solution}{\begin{shaded}\noindent\textbf{Solution.}\ }{\end{shaded}}
\newcommand{\codepart}[1]{\noindent\textit{Programming part. The implementation, documentation and test output are
  in the notebook, section~\textbf{#1}.}\par}
\newcommand{\todo}[1]{\textcolor{red}{\textbf{[TODO: #1]}}}

% hyperref last, as it patches sectioning commands
\usepackage[hidelinks]{hyperref}
